\section{The fundamental group}
Orient every meridian clockwise.  Let $\rho_j$ be the lift of the meridian
about $p_j$ along the zero section of $J$.  Since $B^\circ$ is a pair of
pants,
\[
 \pi_1(J)=\Lambda\rtimes F(\rho_1,\rho_2),
 \qquad
 \rho_j\lambda\rho_j^{-1}=A_j\lambda,
 \qquad
 \rho_0=(\rho_1\rho_2)^{-1}.
\]

\begin{lemma}[Relations imposed by the fillings]
\label{lem:filling-relations}
Under the gluing maps, the clockwise meridian lifts satisfy
\begin{equation}\label{eq:finite-pi1-relations}
 \rho_1^3=t_{v_1},
 \qquad
 \rho_2^4=t_{v_2},
\end{equation}
and the cusp filling kills $\Lambda_{\mathrm{tor}}$ and imposes
\begin{equation}\label{eq:cusp-pi1-relations}
 \rho_1\rho_2=1.
\end{equation}
\end{lemma}

\begin{proof}
By \cref{prop:pi1-Nj}, attaching $N_j$ kills
$\wh g_j^{m_j}t_{-v_j}$.  A clockwise meridian in the $t_j$-disc lifts
through $t_j=s_j^{m_j}$ to a path on which $s_j$ turns through
$-2\pi/m_j$.  Its endpoint is obtained by the deck transformation
$g_j$, because $s_j\circ g_j=e^{-2\pi i/m_j}s_j$.  In the coordinates
of $N_j$, the lift of the zero section is
$\varsigma_j=-\log(s_j)\Pi v_j/(2\pi i)$.  The identity
\[
 R_j(z)\varsigma_j(z)+\frac1{m_j}\Pi(g_jz)v_j
 =\varsigma_j(g_jz)
\]
follows from $\log(s_j\circ g_j)=\log s_j-2\pi i/m_j$ and
$A_jv_j=v_j$.  Thus the gluing identifies $\rho_j$ with $\wh g_j$, with
the same sign for $j=1,2$, and
\eqref{eq:finite-pi1-relations} follows.

By \cref{prop:cusp-boundary}, attaching $N_0$ kills the normal closure of
$\Lambda_{\mathrm{tor}}$ and of the toric meridian.  Under
\cref{prop:cusp-identification}, the zero section is the disc
$(x_1,x_2)=(1,1)$, so its clockwise boundary is the toric meridian and is
also the lift $\rho_0$.  Since $\rho_0=(\rho_1\rho_2)^{-1}$,
\eqref{eq:cusp-pi1-relations} follows.
\end{proof}

\begin{theorem}\label{thm:pi1}
The compact complex threefold $X$ is simply connected.
\end{theorem}

\begin{proof}
Apply van Kampen's theorem to the open cover by the images of $J$ and
slightly thickened cusp and logarithmic pieces.  Their intersections are
connected torus bundles over annuli, so \cref{lem:filling-relations}
gives the quotient of $\pi_1(J)$ by exactly the stated normal relations.
The normal closure of $\Lambda_{\mathrm{tor}}$ under $A_1,A_2$ is
$\ker\gamma$ by \cref{prop:linear-algebra}.  The surviving image of
$\Lambda$ is therefore cyclic, generated by the image $c$ of
$\wh\gamma$, and is central because $A_1,A_2$ act trivially on the
coinvariant quotient.  By \cref{lem:filling-relations}, with
$x=\rho_1$ and $y=\rho_2$, the group has the presentation
\[
 \langle c,x,y\mid c\text{ central},\ xy=1,\ x^3=c,\ y^4=c^{-1}\rangle,
\]
because $\gamma(v_1)=1$ and $\gamma(v_2)=-1$.  Hence $y=x^{-1}$ and
$c=x^3$, while $y^4=c^{-1}$ gives $x^{-4}=x^{-3}$.  Thus $x=1$, and then
$y=c=1$.
\end{proof}

